\scanpage{135}\mspage{470}
leur fils.

Pour le dire autrement encore: à aucun moment en ces premières cinq années
de ma vie, je n'ai connu le sentiment de\quad h o n t e\quad d'\quad ê t r e\quad c e\quad q u e\quad
j e\quad s u i s\quad, que ce soit dans mon corps et ses fonctions, ou dans mes \unsure{pulsion|pulsions}
mes penchants,\unsure{.}mes actions. A aucun moment je n'ai eu à renier quelque chose en
moi, pour être accepté par mon entourage et pouvoir vivre en paix avec lui.

Il arrivait bien sûr que je fasse des choses qui ne ``passaient'' pas: comme
tous les enfants il m'arrivait sûrement d'être pénible, voire insupportable
quand je m'y \unsure{méttais|mettais} - et il était clair parfois qu'il me fallait rectifier le
tir. Je ne faisais pas la loi, ni n'étais tenté de vouloir la faire, n'ayant pas
à compenser quelque mutilation secrète. Et dans l'amour de mes parents pour moi,
il n'aurait pu y avoir de place pour l'adulation, la complaisance aux caprices -
pour une approbation \unsure{inconditionnelle}. Mais s'il arrivait forcément que je me
fasse ``envoyer sur les \unsure{roses|coses}'' par mon père ou par ma mère (tout comme l'inverse
pouvait parfois se produire), jamais dans ces années-là l'un ni l'autre ne m'ont
fait honte, d'un acte ou d'un comportement qui n'aurait pas eu l'heur de leur
plaire.

Sur le fond d'une identification profonde au père, sans \unsure{ambiguité|ambiguïté} aucune,
ma personne comme enfant m'apparaît aujourd'hui comme empreinte à la fois de vi-\newline
rilité et de féminité, fortes l'une comme l'autre.

Il me semble qu'en chaque être et en chaque chose, dans ces indissolubles
et fluctuantes épousailles des qualités yin et yang en lui qui font de lui ce
qu'il est, et dont le délicat équilibre est la beauté profonde, l'harmonie qui
vit en cet être ou en cette chose - que dans cette union intime du yin et du
yang il y a souvent (peut-être toujours) une note de fond, une ``dominante'',
qui est soit yin, soit yang. Cette note de fond n'est pas toujours aisée à dé-\newline
celer en une personne, à cause des mécanismes de répression plus ou moins effi-\newline
caces et complets, qui faussent le jeu en substituant, à une harmonie originelle
une image d'emprunt. Ainsi mon ``image de marque'' pendant quarante ans a été une
image presque exclusivement virile - sans d'ailleurs jamais se voir mettre en
cause ni même être décelée comme telle, par moi-même ni (il me semble) par au-\newline
trui, jusqu'à ma quarante-huitième année. J'ai tendance à croire pourtant que
la note de fond présente à la naissance reste présente pendant la vie entière,
dans des couches profondes tout au moins qui jamais peut-être ne trouveront
l'occasion de se manifester au grand jour. Dans mon propre cas, chose étrange, \unsure{j|je}
ne saurais aujourd'hui encore dire quelle est cette note dominante, celle donc
\mspage{134}
